\scanpage{54}
by $E_{\bar K/K}$) by taking its centralizer $Z(L_i)$ in $E$, whence
\[
(2)\qquad
\diagram{\begin{tikzcd}[ampersand replacement=\&,column sep=small]
1\arrow[r,no head]\&L_i\arrow[r,no head]\arrow[d,phantom,"\simeq" description]\&Z(L_i)\arrow[r]\&\Gamma\arrow[r]\&1\\
\&T_\infty(\bar K^*)=T\&\&\&
\end{tikzcd}}
\]
and we take the splittings of this extension. There are some, for example defined by a uniformizer of $\unsure{\mathcal O}_{\widehat U,i}$. The set of conjugacy classes of the \deleted{\unsure{its}} splittings \deleted{\unsure{the}} of (\unsure{1|2}) is a torsor under \deleted{$H^1(\Gamma,T)$}
\[
(3)\qquad H^1(\Gamma,T)\simeq\varprojlim_n H^1(K,{}_n\mu)\simeq\widehat{K^*}=\varprojlim_n K^*_n
\]
and \deleted{\unsure{this set}} maps injectively \deleted{\unsure{to}} into the set of $\pi$-conjugacy classes of splittings.

\uline{Proposition} We assume $(g,\nu)\ne(0,0),(0,1)$ that is $\Pi_{\bar U}=\Pi_{U/K}\ne(1)$. Then the \inserted{$\pi$-conjugacy} classes of splittings \deleted{\unsure{associated}} of (1) defined by the splittings of (2) are distinct from those associated with the points of $U(K)$. \deleted{Moreover,} If \inserted{(moreover)} $(g,\nu)\ne(0,2)$, \deleted{\unsure{in the}} \inserted{that is if we are} in the anabelian case, then the $\pi$-conjugacy classes of splittings of (1), associated with splittings of (2) for two distinct indices $i\deleted{\unsure{a}}=i_1$ and $i=i_2$ \inserted{are distinct.}
\margin{This is proved except for type $(0,3)$, which deserves a thorough examination.}
\margin{53}
